\par\Needspace{6\baselineskip}\medskip
\begin{center}\hypertarget{native-179}{}\hypertarget{native-179-article132}{}\pdfbookmark[2]{132.}{article-132}132.\end{center}
Der Kürze halber setzen wir
\[
\begin{gathered}
(0.1.2)=a,\qquad (0.\mathrm{I}.2)D'=-b,\\
(0.\mathrm{O}.2)D=c,\qquad (0.\mathrm{II}.2)D''=d,
\end{gathered}
\]
sodass Gleichung 10 in Artikel 114 die Form
\[
a\delta'=b+c\cdot\frac{n}{n'}+d\cdot\frac{n''}{n'}\mathpunct{\raisebox{.45ex}{,}}
\]
annimmt. Die Koeffizienten $c$ und $d$ sind von erster Ordnung. Für die Differenz $c-d$ lässt sich
leicht zeigen, dass sie von zweiter Ordnung ist. Daraus folgt, dass der Wert der Größe $\dfrac{cn+dn''}{n+n''}$, der sich
aus der näherungsweisen Annahme $n:n''=\theta:\theta''$ ergibt, nur mit einem Fehler vierter Ordnung behaftet ist, ja
sogar nur mit einem Fehler fünfter Ordnung, wenn die mittlere Beobachtung von den beiden äußeren
jeweils denselben zeitlichen Abstand hat. Dieser Fehler ist
\[
=\frac{c\theta+d\theta''}{\theta+\theta''\mathpunct{\mkern-8mu\raisebox{-.65ex}{.}\mkern8mu}}
-\frac{cn+dn''}{n+n''}
=\frac{\theta\theta''(d-c)(\eta''-\eta)}{(\theta+\theta'')(\eta''\theta+\eta\theta'')}\mathpunct{\raisebox{.45ex}{,}}
\]
wobei der Nenner von zweiter Ordnung ist. Der eine Faktor $\theta\theta''(d-c)$ des Zählers ist von vierter Ordnung
und der andere, $\eta''-\eta$, von zweiter oder im genannten Sonderfall von dritter Ordnung. Schreibt man die
Gleichung also in der Form
\[
a\delta'=b+\frac{cn+dn''}{n+n''}\cdot\frac{n+n''}{n'}\mathpunct{\raisebox{.45ex}{,}}
\]
so sieht man unmittelbar, dass bei der angenommenen Proportionalität der Größen $n,n''$ zu $\theta,\theta''$
der Mangel des im vorigen Artikel vorgeschlagenen Verfahrens erst durch die \textit{zusätzliche} Annahme entsteht,
dass $n'$ zu $\theta'$ proportional ist. Denn damit wird für den
Faktor $\dfrac{n+n''}{n'}$ der weniger genaue Wert $\dfrac{\theta+\theta''}{\theta'}=1$ eingesetzt, von dem der wahre Wert
\[
=1+\frac{\theta\theta''}{2\eta\eta''rr'r''\cos f\cos f'\cos f''}
\]
um eine Größe zweiter Ordnung abweicht (Artikel 128).
